% GENERATED FILE. Edit canonical lessons, then run npm run generate:books.
% canonical-source-hash: fnv1a64:36636abc1cd97c64
% canonical-lessons: FA-C206-madrak, FA-C206-shenasname, FA-C206-parvande, FA-C206-printer, FA-C206-adres

\chapter{Document, Birth certificate, File, Printer, Address}
\label{ch:fa-document-birth-certificate-file-printer-address}
\hlchaptermodality{206}

\begin{chapteropening}
\textbf{By the end of this chapter:} \emph{I can say a document, a birth certificate, a file, a printer and an address.}

Complete the last lesson of chapter 206: I can say a document, a birth certificate, a file, a printer and an address.
\end{chapteropening}

\section[\texorpdfstring{\fa{مدرک} (madrak)}{مدرک (madrak)}]{\fa{مدرک} (madrak) — a document; a qualification}

\label{lesson:FA-C206-madrak}



\begin{quote}
Before the new one: say the Persian for qeymeh, then the Persian for an olive.
\end{quote}

\subsection*{You'll want to know: \fa{مدرک}}
\textbf{\fa{مدرک}} — \emph{madrak} — \textquotedblleft{}a document; a qualification\textquotedblright{}.

\begin{cousinweb}
\emph{madrak} is borrowed from Arabic \emph{madrak}, \textquotedblleft{}means of grasping, evidence\textquotedblright{}. Its Arabic root is \emph{d-r-k} (Lane). Persian uses it for a document or a qualification.
\end{cousinweb}

\begin{grammarlens}[title={What you've built}]
One more word for this chapter.
\end{grammarlens}

\subsection*{Your turn}
\begin{itemize}
\raggedright
  \item \emph{Say it:} \emph{madrak}
  \item \emph{Say it:} \emph{madrak}, once more
  \item \emph{Say it:} say \emph{madrak}
  \item \emph{Recall:} say the Persian for qeymeh, then the Persian for an olive, then say \emph{madrak} again
\end{itemize}

\begin{tcolorbox}[breakable,colback=teal!4,colframe=teal!35!black,title={Before you move on}]
What does \fa{مدرک} mean? (\textquotedblleft{}A document; a qualification\textquotedblright{}.) Say it once more.
\end{tcolorbox}
\section[\texorpdfstring{\fa{شناسنامه} (shenâsnâme)}{شناسنامه (shenâsnâme)}]{\fa{شناسنامه} (shenâsnâme) — a birth certificate, an ID booklet}

\label{lesson:FA-C206-shenasname}



\begin{quote}
Before the new one: say the Persian for an olive, then the Persian for a document.
\end{quote}

\subsection*{You'll want to know: \fa{شناسنامه}}
\textbf{\fa{شناسنامه}} — \emph{shenâsnâme} — \textquotedblleft{}a birth certificate, an ID booklet\textquotedblright{}.

\begin{cousinweb}
\emph{shenâsnâme} is built from \emph{shenâs}, \textquotedblleft{}knowing, recognising\textquotedblright{} (from \emph{shenâkhtan}), and \emph{nâme}, \textquotedblleft{}letter, book\textquotedblright{}: the book that makes you known. \emph{nâme} is inherited, usually traced to Middle Persian \emph{nāmag}.
\end{cousinweb}

\begin{grammarlens}[title={What you've built}]
One more word for this chapter.
\end{grammarlens}

\subsection*{Your turn}
\begin{itemize}
\raggedright
  \item \emph{Say it:} \emph{shenâsnâme}
  \item \emph{Say it:} \emph{shenâsnâme}, once more
  \item \emph{Say it:} say \emph{shenâsnâme}
  \item \emph{Recall:} say the Persian for an olive, then the Persian for a document, then say \emph{shenâsnâme} again
\end{itemize}

\begin{tcolorbox}[breakable,colback=teal!4,colframe=teal!35!black,title={Before you move on}]
What does \fa{شناسنامه} mean? (\textquotedblleft{}A birth certificate, an ID booklet\textquotedblright{}.) Say it once more.
\end{tcolorbox}
\section[\texorpdfstring{\fa{پرونده} (parvande)}{پرونده (parvande)}]{\fa{پرونده} (parvande) — a file, a dossier}

\label{lesson:FA-C206-parvande}



\begin{quote}
Before the new one: say the Persian for a document, then the Persian for a birth certificate.
\end{quote}

\subsection*{You'll want to know: \fa{پرونده}}
\textbf{\fa{پرونده}} — \emph{parvande} — \textquotedblleft{}a file, a dossier\textquotedblright{}.

\begin{cousinweb}
\emph{parvande} is a native Persian word for a file or bundle, but no source used here gives its older form, and its origin is not settled.
\end{cousinweb}

\begin{grammarlens}[title={What you've built}]
One more word for this chapter.
\end{grammarlens}

\subsection*{Your turn}
\begin{itemize}
\raggedright
  \item \emph{Say it:} \emph{parvande}
  \item \emph{Say it:} \emph{parvande}, once more
  \item \emph{Say it:} say \emph{parvande}
  \item \emph{Recall:} say the Persian for a document, then the Persian for a birth certificate, then say \emph{parvande} again
\end{itemize}

\begin{tcolorbox}[breakable,colback=teal!4,colframe=teal!35!black,title={Before you move on}]
What does \fa{پرونده} mean? (\textquotedblleft{}A file, a dossier\textquotedblright{}.) Say it once more.
\end{tcolorbox}
\section[\texorpdfstring{\fa{پرینتر} (printer)}{پرینتر (printer)}]{\fa{پرینتر} (printer) — a printer}

\label{lesson:FA-C206-printer}



\begin{quote}
Before the new one: say the Persian for a birth certificate, then the Persian for a file.
\end{quote}

\subsection*{You'll want to know: \fa{پرینتر}}
\textbf{\fa{پرینتر}} — \emph{printer} — \textquotedblleft{}a printer\textquotedblright{}.

\begin{cousinweb}
\emph{printer} is a loan, usually traced to English \emph{printer}, \textquotedblleft{}printer\textquotedblright{}.
\end{cousinweb}

\begin{grammarlens}[title={What you've built}]
One more word for this chapter.
\end{grammarlens}

\subsection*{Your turn}
\begin{itemize}
\raggedright
  \item \emph{Say it:} \emph{printer}
  \item \emph{Say it:} \emph{printer}, once more
  \item \emph{Say it:} say \emph{printer}
  \item \emph{Recall:} say the Persian for a birth certificate, then the Persian for a file, then say \emph{printer} again
\end{itemize}

\begin{tcolorbox}[breakable,colback=teal!4,colframe=teal!35!black,title={Before you move on}]
What does \fa{پرینتر} mean? (\textquotedblleft{}A printer\textquotedblright{}.) Say it once more.
\end{tcolorbox}
\section[\texorpdfstring{\fa{آدرس} (âdres)}{آدرس (âdres)}]{\fa{آدرس} (âdres) — an address}

\label{lesson:FA-C206-adres}



\begin{quote}
Before the new one: say the Persian for a file, then the Persian for a printer.
\end{quote}

\subsection*{You'll want to know: \fa{آدرس}}
\textbf{\fa{آدرس}} — \emph{âdres} — \textquotedblleft{}an address\textquotedblright{}.

\begin{cousinweb}
\emph{âdres} is borrowed from French \emph{adresse}, \textquotedblleft{}address\textquotedblright{} (EtymDB).
\end{cousinweb}

\begin{grammarlens}[title={What you've built}]
One more word for this chapter.
\end{grammarlens}

\subsection*{Your turn}
\begin{itemize}
\raggedright
  \item \emph{Say it:} \emph{âdres}
  \item \emph{Say it:} \emph{âdres}, once more
  \item \emph{Say it:} say \emph{âdres}
  \item \emph{Recall:} say the Persian for a file, then the Persian for a printer, then say \emph{âdres} again
\end{itemize}

\begin{tcolorbox}[breakable,colback=teal!4,colframe=teal!35!black,title={Before you move on}]
What does \fa{آدرس} mean? (\textquotedblleft{}An address\textquotedblright{}.) Say it once more.
\end{tcolorbox}
